% Fictitious CV for benchmarks and demos. The person, employers and achievements are invented.
\documentclass[a4paper,11pt]{article}
\usepackage[spanish]{babel}
\usepackage{enumitem}
\usepackage{tabularx}
\newcommand{\resumeItem}[1]{\item\small{#1}}
\newcommand{\resumeSubheading}[4]{\item\begin{tabular*}{0.97\textwidth}[t]{l@{\extracolsep{\fill}}r}\textbf{#1} & #2 \\ \textit{\small#3} & \textit{\small #4} \\\end{tabular*}}
\newcommand{\resumeProjectHeading}[2]{\item\begin{tabular*}{0.97\textwidth}{l@{\extracolsep{\fill}}r}\small#1 & #2 \\\end{tabular*}}
\begin{document}

\begin{center}
  \textbf{\Huge Valentina Soto Pizarro} \\
  \textbf{\large Ingeniera Civil Electrónica} \\
  \small Valparaíso, Chile $|$ valentina.soto@example.com $|$ github.com/example-vsoto
\end{center}

\section{Resumen Profesional}
Ingeniera Civil Electrónica con experiencia en sistemas embebidos, adquisición de datos y
automatización de pruebas. Desarrollo de firmware en \textbf{C/C++} para \textbf{ESP32} y
\textbf{STM32}, integración de sensores y comunicaciones \textbf{MQTT}, y análisis de datos
en \textbf{Python}. Uso de \textbf{Docker} y \textbf{Git} para entornos reproducibles.

\section{Experiencia Profesional}
\begin{itemize}[leftmargin=0.15in, label={}]
  \resumeSubheading{Ingeniera de Desarrollo - Laboratorio de Sensores}{Marzo 2024 -- Actualidad}
    {Instituto Costero de Tecnología (ficticio)}{Valparaíso, Chile}
  \begin{itemize}
    \resumeItem{Desarrollo de firmware en \textbf{C/C++} con \textbf{FreeRTOS} para nodos \textbf{ESP32} de monitoreo ambiental.}
    \resumeItem{Integración de sensores por \textbf{I2C} y \textbf{SPI}; calibración y validación con osciloscopio y analizador lógico.}
    \resumeItem{Publicación de telemetría vía \textbf{MQTT} hacia un servidor Linux; paneles en \textbf{Grafana}.}
    \resumeItem{Scripts de análisis en \textbf{Python} (pandas, NumPy) para detectar derivas de sensores.}
    \resumeItem{Diseño de dos tarjetas \textbf{PCB} de 2 capas en \textbf{KiCad} para acondicionamiento de señales.}
  \end{itemize}
  \resumeSubheading{Práctica Profesional - Mantenimiento Eléctrico}{Enero 2023 -- Marzo 2023}
    {Puerto Interior S.A. (ficticio)}{San Antonio, Chile}
  \begin{itemize}
    \resumeItem{Apoyo en levantamiento de tableros eléctricos y actualización de diagramas unilineales.}
    \resumeItem{Registro de mediciones de vibración en motores de grúas para el plan de mantenimiento.}
  \end{itemize}
\end{itemize}

\section{Proyectos}
\begin{itemize}[leftmargin=0.15in, label={}]
  \resumeProjectHeading{\textbf{Estación meteorológica IoT} $|$ \emph{ESP32, MQTT, Docker, Grafana}}{2024}
  \begin{itemize}
    \resumeItem{Nodo solar con ESP32 que publica datos cada 5 minutos; stack Docker con broker MQTT y Grafana.}
  \end{itemize}
  \resumeProjectHeading{\textbf{Clasificador de defectos en PCB} $|$ \emph{Python, OpenCV}}{2023}
  \begin{itemize}
    \resumeItem{Detección de soldaduras defectuosas con umbralización y contornos en \textbf{OpenCV} (proyecto de curso).}
  \end{itemize}
\end{itemize}

\section{Educación}
\begin{itemize}[leftmargin=0.15in, label={}]
  \resumeSubheading{Universidad del Litoral (ficticia)}{Valparaíso, Chile}
    {Ingeniería Civil Electrónica}{Diciembre 2023}
\end{itemize}

\section{Habilidades Técnicas}
\begin{itemize}[leftmargin=0.15in, label={}]
  \small{\item{
    \textbf{Embebidos}{: ESP32, STM32, FreeRTOS, I2C, SPI, UART, KiCad.} \\
    \textbf{Software}{: C/C++, Python, Git, Docker, Linux.} \\
    \textbf{Datos e IoT}{: MQTT, Grafana, pandas, OpenCV.} \\
    \textbf{Idiomas}{: Español (nativo), Inglés (intermedio).}
  }}
\end{itemize}
\end{document}
